\chapter{Le Quattro Variabili Strutturali in Quadro
Coerente}\label{le-quattro-variabili-strutturali-in-quadro-coerente}

\begin{quote}
\textbf{Argomento portante:} \(\Phi\), \(A\), \(\Omega\),
\(\varepsilon_{\text{dis}}\) non sono parametri indipendenti che entrano
in \(T_s^{(2),\text{sys}}\) come ingredienti di una ricetta --- sono
\textbf{componenti di uno spazio strutturale quadridimensionale
\(\mathbb{V}^{(4)}\)} in cui ogni sistema sociale occupa un punto, e in
cui esistono \emph{sei relazioni di accoppiamento bivariate} che
vincolano la dinamica congiunta. La superficie critica
\(T_s^{(2),\text{sys}} = 1\) è un'ipersurface in \(\mathbb{V}^{(4)}\)
che separa due regimi qualitativamente diversi; le predizioni della
Parte III diventano ora asserzioni sulla \emph{traiettoria} dei quattro
attori in \(\mathbb{V}^{(4)}\) rispetto a tale ipersurface.
\end{quote}



\section{\texorpdfstring{Lo Spazio Strutturale
\(\mathbb{V}^{(4)}\)}{Lo Spazio Strutturale \textbackslash mathbb\{V\}\^{}\{(4)\}}}\label{lo-spazio-strutturale-mathbbv4}

\Hypoth{} \textbf{Definizione.} Sia
\(\mathbb{V}^{(4)} := [0, 1] \times \mathbb{R} \times \mathbb{R}_{\geq 0} \times [0, 1]\)
lo spazio strutturale DEQ$^2$, parametrizzato da:

{\def\LTcaptype{none} % do not increment counter
{\small\begin{longtable}[]{@{}
  >{\raggedright\arraybackslash}p{(\linewidth - 6\tabcolsep) * \real{0.2500}}
  >{\raggedright\arraybackslash}p{(\linewidth - 6\tabcolsep) * \real{0.2500}}
  >{\raggedright\arraybackslash}p{(\linewidth - 6\tabcolsep) * \real{0.2500}}
  >{\raggedright\arraybackslash}p{(\linewidth - 6\tabcolsep) * \real{0.2500}}@{}}
\toprule\noalign{}
\begin{minipage}[b]{\linewidth}\raggedright
Coordinata
\end{minipage} & \begin{minipage}[b]{\linewidth}\raggedright
Variabile
\end{minipage} & \begin{minipage}[b]{\linewidth}\raggedright
Dominio
\end{minipage} & \begin{minipage}[b]{\linewidth}\raggedright
Origine teorica
\end{minipage} \\
\midrule\noalign{}
\endhead
\bottomrule\noalign{}
\endlastfoot
\(\Phi\) & parametro d'ordine di Kuramoto & \([0, 1]\) & Cap. 4 ---
sincronizzazione di fase del campo \(\Pi\) \\
\(A\) & antifragilità & \(\mathbb{R}\) (sub-dominio operativo
\([-1, +\infty)\)) & Cap. 3.3.1 --- \(\partial^2\Pi/\partial\sigma^2\)
(Taleb formalizzato) \\
\(\Omega\) & involuzione & \(\mathbb{R}_{\geq 0}\) & Cap. 3.3.2 ---
\(\Delta E_{\text{diss}}/\Delta E_{\text{utili}} - 1\) (Xiang Biao
formalizzato) \\
\(\varepsilon_{\text{dis}}\) & disimpegno morale & \([0, 1]\) & Cap. 3.2
--- Bandura/Milgram \\
\end{longtable}}
}

\Proved{} \textbf{Conseguenza formale.} Ogni sistema sociale è
descrivibile, ad un dato istante \(t\), come il punto

\[
\boldsymbol{p}(t) = (\Phi(t),\ A(t),\ \Omega(t),\ \varepsilon_{\text{dis}}(t)) \in \mathbb{V}^{(4)}
\]

La sua \emph{traiettoria} è la curva \(t \mapsto \boldsymbol{p}(t)\). La
sua \emph{stabilità sistemica} è il valore scalare
\(T_s^{(2),\text{sys}}(\boldsymbol{p}(t); s_{iii}, z, \delta, \sigma)\),
dove i parametri ereditati dalla Geometria/Topografia
\((s_{iii}, z, \delta, \sigma)\) sono trattati come \emph{campi esogeni}
(lentamente variabili rispetto alla scala temporale di
\(\boldsymbol{p}\)).

\Hypoth{} \textbf{Riduzione dimensionale operativa.}
\(\mathbb{V}^{(4)}\) ammette una proiezione canonica nei due piani
notevoli:

\begin{itemize}
\tightlist
\item
  \(\mathbb{V}^{(4)} \to \mathbb{V}^{(2)}_{\text{coerenza}} := (\Phi, A)\)
  --- \textbf{piano della coerenza-valore};
\item
  \(\mathbb{V}^{(4)} \to \mathbb{V}^{(2)}_{\text{decoerenza}} := (\Omega, \varepsilon_{\text{dis}})\)
  --- \textbf{piano della decoerenza-spreco}.
\end{itemize}

Le ipersurfaci di \(T_s^{(2),\text{sys}}\) sono \emph{prodotti} di curve
in questi due piani --- proprietà che useremo ripetutamente nel Cap. 10
per le dashboard SPC-DEQ.



\section{Le Sei Relazioni di Accoppiamento
Bivariate}\label{le-sei-relazioni-di-accoppiamento-bivariate}

\Hypoth{} \textbf{Tesi.} Le quattro variabili \emph{non sono
indipendenti}. Esistono \(\binom{4}{2} = 6\) relazioni di accoppiamento
bivariate, ciascuna con un segno e una giustificazione strutturale. Esse
vincolano lo spazio dei punti effettivamente realizzabili in
\(\mathbb{V}^{(4)}\) --- chiameremo questo sottoinsieme
\(\mathbb{V}^{(4)}_{\text{ammissibile}}\).

{\def\LTcaptype{none} % do not increment counter
{\small\begin{longtable}[]{@{}
  >{\raggedright\arraybackslash}p{(\linewidth - 8\tabcolsep) * \real{0.2000}}
  >{\raggedright\arraybackslash}p{(\linewidth - 8\tabcolsep) * \real{0.2000}}
  >{\raggedright\arraybackslash}p{(\linewidth - 8\tabcolsep) * \real{0.2000}}
  >{\raggedright\arraybackslash}p{(\linewidth - 8\tabcolsep) * \real{0.2000}}
  >{\raggedright\arraybackslash}p{(\linewidth - 8\tabcolsep) * \real{0.2000}}@{}}
\toprule\noalign{}
\begin{minipage}[b]{\linewidth}\raggedright
\#
\end{minipage} & \begin{minipage}[b]{\linewidth}\raggedright
Coppia
\end{minipage} & \begin{minipage}[b]{\linewidth}\raggedright
Segno
\end{minipage} & \begin{minipage}[b]{\linewidth}\raggedright
Meccanismo
\end{minipage} & \begin{minipage}[b]{\linewidth}\raggedright
Verifica nei Cap. 6-9
\end{minipage} \\
\midrule\noalign{}
\endhead
\bottomrule\noalign{}
\endlastfoot
\textbf{R1} & \(\Phi \leftrightarrow A\) & \(+\) & Coerenza permette
redistribuzione del guadagno antifragile & Brasile: \(\Phi\) moderato
\(+\) \(A\) adattivo \(\Rightarrow A_{\text{sys}} > 0\) \\
\textbf{R2} & \(\Phi \leftrightarrow \Omega\) & \(-\) & Coerenza
inibisce duplicazione e competizione improduttiva & Cina:
\(\Phi^{\text{imposto}}\) alto, \(\Omega^{\text{spontaneo}}\)
esplosivo \\
\textbf{R3} & \(\Phi \leftrightarrow \varepsilon_{\text{dis}}\) & \(-\)
& Coerenza è correlata a \(Q_0\) alto \(\Rightarrow\) basso disimpegno &
USA: \(\Phi \to 0 \Rightarrow\) 40\%+ vede l'altro come \emph{minaccia}
(Cap. 7.1) \\
\textbf{R4} & \(A \leftrightarrow \Omega\) & \(-\) (forte) &
Antifragilità trasforma volatilità in opzionalità \(\Rightarrow\) riduce
dissipazione strutturale & USA: \(A_{\text{sys}} < 0\) (iper-estensione)
\(\Rightarrow \Omega \uparrow\) (conflitto interno) \\
\textbf{R5} & \(A \leftrightarrow \varepsilon_{\text{dis}}\) & \(-\)
(debole) & Sistemi antifragili tollerano dissenso \(\Rightarrow\)
riducono disimpegno morale & Brasile: \(A^{\text{adattivo}}\) alto \(+\)
\(\varepsilon\) in calo post-2022 \\
\textbf{R6} & \(\Omega \leftrightarrow \varepsilon_{\text{dis}}\) &
\(+\) & Involuzione produce frustrazione \(\Rightarrow\) disimpegno
morale crescente & Cina: \(\Omega^{\text{Neijuan}}\) \(\Rightarrow\)
\emph{lying flat}, fenomeno generazionale \\
\end{longtable}}
}

\Proved{} \textbf{Risultato strutturale.} Quattro relazioni hanno segno
\emph{negativo} (R2, R3, R4, R5) e due hanno segno \emph{positivo} (R1,
R6). Strutturalmente: \textbf{valore e coerenza sono coalleate; spreco e
disimpegno sono coalleati}. Sistemi \emph{miscelati} (alta \(\Phi\) con
alto \(\Omega\), oppure alto \(A\) con alto
\(\varepsilon_{\text{dis}}\)) sono possibili ma metastabili ---
appartengono al \emph{bordo} di
\(\mathbb{V}^{(4)}_{\text{ammissibile}}\) e tendono a evolvere verso uno
dei due assi.

\Hypoth{} \textbf{Asimmetria fondamentale.} La forza di R4
(\(A\)-\(\Omega\)) è strutturalmente la più alta perché entrambe le
variabili sono \emph{misurate sullo stesso flusso
energetico-produttivo}: \(A\) ne misura la convessità rispetto allo
shock, \(\Omega\) ne misura la dissipazione interna. Si spiega così
perché il Cap. 7 ha potuto inferire \(A_{\text{sys}}^{\text{USA}} < 0\)
dall'evidenza \(\Omega^{\text{USA}} \uparrow\) senza misurare \(A\)
direttamente: R4 lo permette.



\section{\texorpdfstring{Mappa dei Quattro Attori in
\(\mathbb{V}^{(4)}\)}{Mappa dei Quattro Attori in \textbackslash mathbb\{V\}\^{}\{(4)\}}}\label{mappa-dei-quattro-attori-in-mathbbv4}

\Proved{} \textbf{Posizioni stimate aprile 2026} (su scala normalizzata;
\(\Phi, \varepsilon_{\text{dis}} \in [0,1]\), \(A\) in unità di
\(\sigma_0^{-2}\) con range tipico \([-1, +1]\), \(\Omega\)
adimensionale con range tipico \([0, 2]\)):

{\def\LTcaptype{none} % do not increment counter
{\small\begin{longtable}[]{@{}
  >{\raggedright\arraybackslash}p{(\linewidth - 10\tabcolsep) * \real{0.1667}}
  >{\raggedright\arraybackslash}p{(\linewidth - 10\tabcolsep) * \real{0.1667}}
  >{\raggedright\arraybackslash}p{(\linewidth - 10\tabcolsep) * \real{0.1667}}
  >{\raggedright\arraybackslash}p{(\linewidth - 10\tabcolsep) * \real{0.1667}}
  >{\raggedright\arraybackslash}p{(\linewidth - 10\tabcolsep) * \real{0.1667}}
  >{\raggedright\arraybackslash}p{(\linewidth - 10\tabcolsep) * \real{0.1667}}@{}}
\toprule\noalign{}
\begin{minipage}[b]{\linewidth}\raggedright
Attore
\end{minipage} & \begin{minipage}[b]{\linewidth}\raggedright
\(\Phi\)
\end{minipage} & \begin{minipage}[b]{\linewidth}\raggedright
\(A\)
\end{minipage} & \begin{minipage}[b]{\linewidth}\raggedright
\(\Omega\)
\end{minipage} & \begin{minipage}[b]{\linewidth}\raggedright
\(\varepsilon_{\text{dis}}\)
\end{minipage} & \begin{minipage}[b]{\linewidth}\raggedright
\(T_s^{(2),\text{sys}}\) stimato
\end{minipage} \\
\midrule\noalign{}
\endhead
\bottomrule\noalign{}
\endlastfoot
\textbf{Solitone (Musk)} & indef. (\(N=1\));
\(\Phi^{\text{esogeno}} \approx 0{,}3\) & \(+0{,}7\) individuale &
\(-0{,}1\) (anti-spreco operativo) & \(0{,}5\) (tolleranza alta a
violazione norme) & non applicabile sistemicamente; influenza esogena
\(\zeta(t)\) \\
\textbf{Egemone Decoerente (USA)} & \(0{,}25\)-\(0{,}35\) (V-Dem
\(0{,}57\), Pew, Brand Finance) & \(-0{,}3\) (sistemico, da
iper-estensione) & \(0{,}8\)-\(1{,}2\) (polarizzazione, tariffe
auto-shock) & \(0{,}55\)-\(0{,}65\) (40\%+ vede l'altro come minaccia) &
\textbf{\(\approx 0{,}45\)}, sotto soglia \\
\textbf{Rigido Sincronizzato (Cina)} &
\(\Phi^{\text{oss}} \approx 0{,}90\),
\(\Phi^{\text{spont}} \approx 0{,}40\)-\(0{,}50\) &
\(A^{\text{tecnologico}} +0{,}5\), \(A^{\text{politico}} -0{,}4\) &
\(1{,}0\)-\(1{,}4\) (Neijuan documentato + property + deflazione 10
trim.) & \(0{,}40\)-\(0{,}50\) (lying flat, anti-996, alienazione
generazionale) & \textbf{\(\approx 0{,}80\) apparente,
\(\approx 0{,}50\) effettivo} \\
\textbf{Bilanciere (Brasile)} & \(0{,}55\)-\(0{,}65\) (V-Dem U-turn,
\(R_c\) effettivo) & \(+0{,}4\) (resilienza 2018-2023 confermata) &
\(0{,}3\)-\(0{,}5\) (fiscale tirato, ma non involutivo) &
\(0{,}40\)-\(0{,}55\) (polarizzazione persistente, ma sub-USA) &
\textbf{\(\approx 1{,}10\)}, sopra soglia condizionale \\
\end{longtable}}
}

\Hypoth{} \textbf{Lettura strutturale.} Solo il Brasile ha il punto
\(\boldsymbol{p}(t)\) \textbf{interno} alla regione
\(T_s^{(2),\text{sys}} > 1\) in modo non degenere. USA, Cina e Solitone
sono in tre forme distinte di patologia strutturale:

\begin{itemize}
\tightlist
\item
  USA: punto \emph{interno} alla regione \(T_s^{(2),\text{sys}} < 1\)
  (decoerenza distribuita stabile);
\item
  Cina: punto \emph{sul bordo} della regione critica, mantenuto dalla
  forzatura \(K^{\text{imposto}}\) (instabile a perturbazioni di \(K\));
\item
  Solitone: punto \emph{fuori dominio} sistemico (\(N=1\), \(\Phi\)
  degenere).
\end{itemize}

\Proved{} \textbf{Verifica indiretta della consistenza.} Le posizioni
\(\boldsymbol{p}_{\text{USA}}\), \(\boldsymbol{p}_{\text{Cina}}\),
\(\boldsymbol{p}_{\text{Brasile}}\) sono consistenti con tutte e sei le
relazioni R1-R6 simultaneamente --- il che non era garantito a priori e
costituisce una validazione interna dell'apparato.



\section{\texorpdfstring{La Superficie Critica
\(T_s^{(2),\text{sys}} = 1\)}{La Superficie Critica T\_s\^{}\{(2),\textbackslash text\{sys\}\} = 1}}\label{la-superficie-critica-t_s2textsys-1}

\Hypoth{} \textbf{Costruzione.} L'equazione

\[
T_s^{(2),\text{sys}}(\Phi, A, \Omega, \varepsilon_{\text{dis}}; s_{iii}, z, \delta, \sigma) = 1
\]

definisce, per ciascuna scelta di \((s_{iii}, z, \delta, \sigma)\)
esogeni, un'\textbf{ipersurface tridimensionale}
\(\mathcal{S} \subset \mathbb{V}^{(4)}\) --- la \emph{frontiera di
transizione di fase}. Esplicitando dalla forma del Cap. 4:

\[
s_{iii} \cdot z \cdot \delta \cdot (1 + \Phi A) = \sigma \cdot \varepsilon_{\text{dis}} \cdot \left(\Phi + \frac{\Omega}{\Phi} \cdot \Phi\right) = \sigma \cdot \varepsilon_{\text{dis}} \cdot (\Phi + \Omega)
\]

\Hypoth{} \textbf{Forma quadratica esatta.} Imponendo
\(T_s^{(2),\text{sys}} = 1\) e moltiplicando per \(\Phi\) entrambi i
lati si ottiene un'equazione quadratica in \(\Phi^*\):

\[
s_{iii} z \delta \langle A\rangle \cdot (\Phi^*)^2 + (s_{iii} z \delta - \sigma \varepsilon_{\text{dis}}) \cdot \Phi^* - \sigma \varepsilon_{\text{dis}} \langle\Omega\rangle = 0
\]

\Hypoth{} \textbf{Approssimazione lineare di servizio} (per
\(\Phi^* \approx 1\), regime di alta coerenza):

\[
\Phi^*_{\text{lin}}(A, \Omega, \varepsilon_{\text{dis}}) \approx \frac{\sigma \varepsilon_{\text{dis}} \Omega - s_{iii} z \delta}{s_{iii} z \delta A - \sigma \varepsilon_{\text{dis}}}
\]

(coerente con Cap. 4.5.1; la \textbf{forma esatta} via formula
quadratica è derivata in \textbf{Appendice A \S{}A.7} ed è preferibile per
sistemi in regime di forte decoerenza, dove \(\Phi\) è lontano da 1).

\Proved{} \textbf{Proprietà della superficie \(\mathcal{S}\).}

{\def\LTcaptype{none} % do not increment counter
{\small\begin{longtable}[]{@{}
  >{\raggedright\arraybackslash}p{(\linewidth - 4\tabcolsep) * \real{0.3333}}
  >{\raggedright\arraybackslash}p{(\linewidth - 4\tabcolsep) * \real{0.3333}}
  >{\raggedright\arraybackslash}p{(\linewidth - 4\tabcolsep) * \real{0.3333}}@{}}
\toprule\noalign{}
\begin{minipage}[b]{\linewidth}\raggedright
Sezione
\end{minipage} & \begin{minipage}[b]{\linewidth}\raggedright
Forma
\end{minipage} & \begin{minipage}[b]{\linewidth}\raggedright
Interpretazione
\end{minipage} \\
\midrule\noalign{}
\endhead
\bottomrule\noalign{}
\endlastfoot
Sezione \(\Omega = 0\) & \(\Phi^* = -1/A\) se \(A < 0\); non definita se
\(A > 0\) & senza involuzione, solo sistemi \emph{fragili} possono
attraversare la soglia \\
Sezione \(A = 0\) &
\(\Phi^* = (\sigma\varepsilon\Omega - s_{iii}z\delta) / (-\sigma\varepsilon)\)
& senza antifragilità, \(\Phi^*\) cresce linearmente con \(\Omega\) \\
Sezione \(\varepsilon_{\text{dis}} = 0\) &
\(\Phi^* = -s_{iii}z\delta / (s_{iii}z\delta A)\) se \(A > 0\), indef.
altrimenti & senza disimpegno morale, sistema tipicamente in regime
sicuro \\
Sezione \(\Phi = 1\) &
\(1 + A = (\sigma\varepsilon/s_{iii}z\delta)(1 + \Omega)\) & recupero
della soglia DEQ$^1$ con \(A, \Omega\) \\
\end{longtable}}
}

\Hypoth{} \textbf{Topologia di \(\mathcal{S}\).} \(\mathcal{S}\) divide
\(\mathbb{V}^{(4)}_{\text{ammissibile}}\) in due regioni:

\begin{itemize}
\tightlist
\item
  \(\mathcal{R}_+ := \{\boldsymbol{p} : T_s^{(2),\text{sys}}(\boldsymbol{p}) > 1\}\)
  --- regione di \textbf{stabilità sistemica};
\item
  \(\mathcal{R}_- := \{\boldsymbol{p} : T_s^{(2),\text{sys}}(\boldsymbol{p}) < 1\}\)
  --- regione di \textbf{transizione di fase imminente} (Cap. 3.5.3,
  \(\Delta t > \Delta t^*\)).
\end{itemize}

\Hypoth{} \textbf{Distanza alla soglia.} Definiamo la \emph{distanza
alla criticità} di un punto \(\boldsymbol{p}\) come
\(d(\boldsymbol{p}, \mathcal{S})\) con metrica sulla scala normalizzata.
Il segno è positivo se \(\boldsymbol{p} \in \mathcal{R}_+\), negativo
altrimenti. Questa è la quantità scalare che la dashboard SPC-DEQ del
Cap. 10 dovrà monitorare.



\section{Dinamica delle Quattro
Variabili}\label{dinamica-delle-quattro-variabili}

\Hypoth{} \textbf{Postulato dinamico.} Le quattro variabili evolvono
secondo un sistema di equazioni differenziali accoppiate, con una
struttura formale ereditata dal Kuramoto generalizzato:

\[
\begin{aligned}
\dot{\Phi}(t) &= K(t) \cdot f_\Phi(\Phi) - \gamma_\Phi \cdot \zeta(t) - \beta_\Phi \cdot \varepsilon_{\text{dis}}(t) \\
\dot{A}(t) &= -\alpha_A \cdot (A - A_\infty(\Phi, \Omega)) + \eta_A(t) \\
\dot{\Omega}(t) &= \alpha_\Omega \cdot \mathcal{F}(\langle\Pi\rangle - \Pi^*) - \mu_\Omega \cdot \Phi \cdot A + \eta_\Omega(t) \\
\dot{\varepsilon}_{\text{dis}}(t) &= \alpha_\varepsilon \cdot \Omega - \mu_\varepsilon \cdot \Phi \cdot Q_0(t) + \eta_\varepsilon(t)
\end{aligned}
\]

dove:

\begin{itemize}
\tightlist
\item
  \(K(t)\) è la forza di accoppiamento di Kuramoto del Cap. 4 (esogena,
  dipende da infrastruttura comunicativa e istituzionale);
\item
  \(\zeta(t)\) è la sorgente esogena di disordine (Cap. 6.6: il Solitone
  come \(\zeta\) tipico per il sistema USA);
\item
  \(A_\infty(\Phi, \Omega)\) è il valore asintotico di \(A\) a
  \((\Phi, \Omega)\) fissati --- relazione R1-R4 incorporata;
\item
  \(\mathcal{F}(\cdot)\) è una funzione monotona crescente che modella
  l'auto-rinforzo dell'involuzione quando il payoff aggregato scende
  sotto il livello di soglia \(\Pi^*\);
\item
  \(Q_0(t)\) è la qualità etica del campo discorsivo (eredità Topografia
  \S{}5.2);
\item
  \(\eta_X(t)\) sono rumori stocastici a media nulla.
\end{itemize}

\Proved{} \textbf{Conseguenza interpretativa.} Le sei relazioni R1-R6 di
\$\S\$5.1 emergono come \emph{equilibri quasi-statici} del sistema
dinamico: i segni dei coefficienti \(\beta, \mu, \alpha\) realizzano
R1-R6 nelle equazioni evolutive. Per esempio,
\(-\beta_\Phi \varepsilon_{\text{dis}}\) realizza R3;
\(-\mu_\Omega \Phi A\) realizza R2 e R4 simultaneamente.

\Conj{} \textbf{Caveat.} Il sistema dinamico nella forma di sopra è un
\emph{postulato operativo}: la sua calibrazione empirica richiede
dataset multi-paese su 20+ anni (rinviata a paper EN). Il valore
strutturale è che fornisce uno \emph{schema di pensiero} per organizzare
gli indicatori monitorabili del Cap. 10.

\Hypoth{} \textbf{Applicazione ai quattro attori (verifica
qualitativa).}

{\def\LTcaptype{none} % do not increment counter
{\small\begin{longtable}[]{@{}
  >{\raggedright\arraybackslash}p{(\linewidth - 4\tabcolsep) * \real{0.3333}}
  >{\raggedright\arraybackslash}p{(\linewidth - 4\tabcolsep) * \real{0.3333}}
  >{\raggedright\arraybackslash}p{(\linewidth - 4\tabcolsep) * \real{0.3333}}@{}}
\toprule\noalign{}
\begin{minipage}[b]{\linewidth}\raggedright
Attore
\end{minipage} & \begin{minipage}[b]{\linewidth}\raggedright
Termine dominante
\end{minipage} & \begin{minipage}[b]{\linewidth}\raggedright
Conseguenza dinamica
\end{minipage} \\
\midrule\noalign{}
\endhead
\bottomrule\noalign{}
\endlastfoot
\textbf{USA} & \(-\gamma_\Phi \zeta(t)\) con
\(\zeta = \zeta_{\text{Sol}}(t)\) & \(\dot\Phi < 0\) persistente
\(\Rightarrow \Phi \to 0\) (Cap. 7.4) \\
\textbf{Cina} & \(K(t)\) massimo da imposizione \(+\)
\(\alpha_\Omega \mathcal{F}\) esplosivo & \(\Phi^{\text{osservato}}\)
stabile, ma \(\dot\Omega > 0\) persistente (Cap. 8.2) \\
\textbf{Brasile} & \(-\mu_\varepsilon \Phi Q_0\) con \(Q_0\) alto
post-2022 & \(\dot\varepsilon < 0\) post-Bolsonaro \(\Rightarrow\)
U-turn V-Dem (Cap. 9.4) \\
\textbf{Solitone} & non applicabile sistemicamente; agisce come
\(\zeta\) esogeno & influenza altri sistemi via accoppiamento esogeno
(Cap. 6.6, 7.4) \\
\end{longtable}}
}



\section{Gradi di Libertà
Operativi}\label{gradi-di-libertuxe0-operativi}

\Hypoth{} \textbf{Domanda strategica.} Quali variabili sono
\emph{modificabili} da scelte politiche, e su quale scala temporale?

{\def\LTcaptype{none} % do not increment counter
{\small\begin{longtable}[]{@{}
  >{\raggedright\arraybackslash}p{(\linewidth - 6\tabcolsep) * \real{0.2500}}
  >{\raggedright\arraybackslash}p{(\linewidth - 6\tabcolsep) * \real{0.2500}}
  >{\raggedright\arraybackslash}p{(\linewidth - 6\tabcolsep) * \real{0.2500}}
  >{\raggedright\arraybackslash}p{(\linewidth - 6\tabcolsep) * \real{0.2500}}@{}}
\toprule\noalign{}
\begin{minipage}[b]{\linewidth}\raggedright
Variabile
\end{minipage} & \begin{minipage}[b]{\linewidth}\raggedright
Modificabilità
\end{minipage} & \begin{minipage}[b]{\linewidth}\raggedright
Scala temporale tipica
\end{minipage} & \begin{minipage}[b]{\linewidth}\raggedright
Leve operative
\end{minipage} \\
\midrule\noalign{}
\endhead
\bottomrule\noalign{}
\endlastfoot
\(\Phi\) & media & 3-7 anni & infrastruttura comunicativa, integrazione
istituzionale, riforma elettorale \\
\(A\) & bassa-media & 5-15 anni & educazione, ridondanza istituzionale,
opzionalità sistemica (Taleb) \\
\(\Omega\) & media-alta & 1-5 anni & regolazione mercati, anti-trust,
\emph{anti-involution} effettiva (vs.~dichiarativa) \\
\(\varepsilon_{\text{dis}}\) & alta sul breve, bassa sul lungo & 6
mesi-3 anni & qualità del discorso pubblico (\(Q_0\)), enforcement
giudiziario, esempi pubblici \\
\end{longtable}}
}

\Proved{} \textbf{Conseguenza strategica della Parte V.} Le strategie
antifragili del Cap. 12 (individuali, organizzative, statali) saranno
strutturate come \emph{interventi su sottospazi specifici di
\(\mathbb{V}^{(4)}\)} a scala temporale appropriata. Non esistono
\emph{silver bullets} --- esistono solo combinazioni coerenti.

\Conj{} \textbf{Limite irriducibile.} Le variabili esogene
\((s_{iii}, z, \delta, \sigma)\) sono parzialmente fuori controllo
(specialmente \(\sigma\) --- volatilità ambientale che include shock
climatici, geopolitici, tecnologici). La DEQ$^2$ non promette
stabilizzazione \emph{deterministica} --- promette \emph{posizionamento
in \(\mathcal{R}_+\) con margine di sicurezza} rispetto a
\(\mathcal{S}\).



\section{Indicatori Osservabili e Protocollo di
Misurazione}\label{indicatori-osservabili-e-protocollo-di-misurazione}

\Proved{} \textbf{Mappa indicatori-variabili} (preparazione Cap. 10):

{\def\LTcaptype{none} % do not increment counter
{\small\begin{longtable}[]{@{}
  >{\raggedright\arraybackslash}p{(\linewidth - 6\tabcolsep) * \real{0.2500}}
  >{\raggedright\arraybackslash}p{(\linewidth - 6\tabcolsep) * \real{0.2500}}
  >{\raggedright\arraybackslash}p{(\linewidth - 6\tabcolsep) * \real{0.2500}}
  >{\raggedright\arraybackslash}p{(\linewidth - 6\tabcolsep) * \real{0.2500}}@{}}
\toprule\noalign{}
\begin{minipage}[b]{\linewidth}\raggedright
Variabile
\end{minipage} & \begin{minipage}[b]{\linewidth}\raggedright
Proxy primario
\end{minipage} & \begin{minipage}[b]{\linewidth}\raggedright
Proxy secondari
\end{minipage} & \begin{minipage}[b]{\linewidth}\raggedright
Frequenza
\end{minipage} \\
\midrule\noalign{}
\endhead
\bottomrule\noalign{}
\endlastfoot
\(\Phi\) & V-Dem Liberal Democracy Index & Pew trust nel governo; Pew
common ground; Polarization Index; partisan ANES & annuale \\
\(A\) & resilienza 2008-2009 + COVID-2020 + 2022-2023 (3 stress tests) &
Total Factor Productivity convexity; opzioni reali aziendali; ridondanza
istituzionale (federalismo) & retrospettiva (ogni stress event) \\
\(\Omega\) & Total Factor Productivity in declino persistente &
property/sector-specific stress tests (Cina); judicial productivity;
deflazione PPI & trimestrale \\
\(\varepsilon_{\text{dis}}\) & Pew ``altro partito come minaccia'';
\emph{disengagement} sondaggi & turnout elettorale; volontariato;
protest participation; trust mediatico & annuale \\
\end{longtable}}
}

\Hypoth{} \textbf{Calibrazione cross-sezionale.} Il dataset minimo per
stimare i coefficienti del sistema dinamico è \(50+\) paesi
\(\times 20+\) anni (1995-2025) --- esattamente il dataset target del
paper EN (Appendice B). I quattro attori del libro sono casi-studio
illustrativi, non base statistica.

\Conj{} \textbf{Limite metodologico.} \(A\) è la variabile più difficile
da misurare in tempo reale --- richiede \emph{eventi di stress} per
mostrarsi. Il Cap. 10 propone l'uso di \emph{simulated stress tests}
(Monte Carlo controfattuali) come sostituto operativo.



\section{Sintesi del Capitolo}\label{sintesi-del-capitolo}

\Proved{} \textbf{In una frase:}
\(\mathbb{V}^{(4)} = (\Phi, A, \Omega, \varepsilon_{\text{dis}})\) è lo
spazio strutturale in cui ogni sistema sociale è un punto, vincolato da
sei relazioni di accoppiamento bivariate (R1-R6) e dotato di una
superficie critica \(\mathcal{S}\) che separa i regimi di stabilità e di
transizione di fase imminente; le strategie politiche del libro saranno
\emph{traiettorie programmate} nel sottospazio
\(\mathbb{V}^{(4)}_{\text{ammissibile}}\) con vincolo
\(d(\boldsymbol{p}, \mathcal{S}) > 0\).

{\def\LTcaptype{none} % do not increment counter
{\small\begin{longtable}[]{@{}
  >{\raggedright\arraybackslash}p{(\linewidth - 4\tabcolsep) * \real{0.3333}}
  >{\raggedright\arraybackslash}p{(\linewidth - 4\tabcolsep) * \real{0.3333}}
  >{\raggedright\arraybackslash}p{(\linewidth - 4\tabcolsep) * \real{0.3333}}@{}}
\toprule\noalign{}
\begin{minipage}[b]{\linewidth}\raggedright
Componente
\end{minipage} & \begin{minipage}[b]{\linewidth}\raggedright
Funzione
\end{minipage} & \begin{minipage}[b]{\linewidth}\raggedright
Capitolo successivo che la usa
\end{minipage} \\
\midrule\noalign{}
\endhead
\bottomrule\noalign{}
\endlastfoot
Spazio \(\mathbb{V}^{(4)}\) + relazioni R1-R6 & Quadro unificato & Cap.
10 dashboard SPC-DEQ \\
Mappa quattro attori & Validazione interna apparato & Cap. 11 predizioni
datate \\
Superficie \(\mathcal{S}\) + distanza \(d(\boldsymbol{p}, \mathcal{S})\)
& Indicatore scalare di stabilità & Cap. 10 dashboard + Cap. 11
falsificazione \\
Sistema dinamico \(\dot\Phi, \dot A, \dot\Omega, \dot\varepsilon\) &
Schema evolutivo & Cap. 10 scenari quantizzati \\
Gradi di libertà operativi & Mappa strategica & Cap. 12 protocolli
antifragili \\
Mappa indicatori-variabili & Operativizzazione & Cap. 10 dashboard \(+\)
Appendice C \\
\end{longtable}}
}

\Hypoth{} \textbf{Tesi conclusiva.} Il Cap. 4 ha derivato
\(T_s^{(2),\text{sys}}\) come \emph{funzione scalare} delle quattro
variabili strutturali. Il Cap. 5 mostra che quelle variabili sono
\emph{intrinsecamente accoppiate} (R1-R6) e abitano uno spazio
\(\mathbb{V}^{(4)}\) con geometria propria (la superficie
\(\mathcal{S}\)). La Parte II del libro è ora chiusa: l'apparato
unificato è completo, la Parte III ne ha verificato l'applicabilità sui
quattro attori paradigmatici, e la Parte IV può procedere alla
\emph{operativizzazione} tramite dashboard, scenari quantizzati e
predizioni rischiose.



\section{Note Epistemiche di Chiusura
Capitolo}\label{note-epistemiche-di-chiusura-capitolo}

{\def\LTcaptype{none} % do not increment counter
{\small\begin{longtable}[]{@{}
  >{\raggedright\arraybackslash}p{(\linewidth - 4\tabcolsep) * \real{0.3333}}
  >{\raggedright\arraybackslash}p{(\linewidth - 4\tabcolsep) * \real{0.3333}}
  >{\raggedright\arraybackslash}p{(\linewidth - 4\tabcolsep) * \real{0.3333}}@{}}
\toprule\noalign{}
\begin{minipage}[b]{\linewidth}\raggedright
\#
\end{minipage} & \begin{minipage}[b]{\linewidth}\raggedright
Affermazione
\end{minipage} & \begin{minipage}[b]{\linewidth}\raggedright
Status
\end{minipage} \\
\midrule\noalign{}
\endhead
\bottomrule\noalign{}
\endlastfoot
1 & Definizione di \(\mathbb{V}^{(4)}\) come spazio strutturale &
\Hypoth{} costruzione formale \\
2 & Sei relazioni R1-R6 con segni e meccanismi & \Hypoth{} quadro
analitico, segni verificati nei Cap. 6-9 \\
3 & Asimmetria notevole: 4 segni negativi, 2 positivi & \Proved{}
struttura termodinamica del sistema sociale \\
4 & R4 (\(A\)-\(\Omega\)) come relazione più forte & \Hypoth{} inferenza
strutturale, usata operativamente nel Cap. 7 \\
5 & Mappa quattro attori in \(\mathbb{V}^{(4)}\) con stime numeriche &
\Conj{} stime parametriche con intervalli di confidenza ampi \\
6 & Superficie critica \(\mathcal{S}\) con forma esplicita
\(\Phi^*(A, \Omega, \varepsilon)\) & \Proved{} derivata analiticamente
da Cap. 4 \\
7 & Topologia di \(\mathbb{V}^{(4)}\) in \(\mathcal{R}_+\) e
\(\mathcal{R}_-\) & \Proved{} conseguenza diretta di
\(T_s^{(2),\text{sys}} = 1\) \\
8 & Sistema dinamico per
\(\dot\Phi, \dot A, \dot\Omega, \dot\varepsilon\) & \Hypoth{} postulato
operativo, calibrazione rinviata a paper EN \\
9 & Mappa gradi di libertà operativi \(\to\) scale temporali \(\to\)
leve & \Hypoth{} sintesi strategica, base per Cap. 12 \\
10 & Mappa indicatori-variabili per dashboard & \Proved{} derivata da
fonti citate nei Cap. 7-9 \\
11 & Validazione interna: posizioni
\(\boldsymbol{p}_{\text{USA, Cina, Brasile}}\) consistenti con tutte
R1-R6 & \Proved{} verifica non banale dell'apparato \\
\end{longtable}}
}



\section{Chiusura della Parte II e Apertura della Parte
IV}\label{chiusura-della-parte-ii-e-apertura-della-parte-iv}

La Parte II del libro è ora completa: Cap. 3 ha derivato \(T_s^{(2)}\)
locale, Cap. 4 ha derivato \(T_s^{(2),\text{sys}}\) con asimmetria
\(\Phi\)-moltiplica-\(A\) vs \(\Phi\)-divide-\(\Omega\), Cap. 5 ha
consolidato le quattro variabili come spazio strutturale
\(\mathbb{V}^{(4)}\) con sei relazioni di accoppiamento. La Parte III ha
mappato i quattro attori paradigmatici nello spazio \(\mathbb{V}^{(4)}\)
(Cap. 6-9). La \textbf{Parte IV} può ora trasformare l'apparato in
\textbf{strumento operativo}:

\begin{itemize}
\tightlist
\item
  \textbf{Cap. 10} --- Scenari quantizzati e Dashboard SPC-DEQ: le
  proiezioni \(\mathbb{V}^{(4)} \to \mathbb{V}^{(2)}_{\text{coerenza}}\)
  e \(\mathbb{V}^{(4)} \to \mathbb{V}^{(2)}_{\text{decoerenza}}\)
  diventano cruscotti di monitoraggio in tempo reale, con soglie
  operative e segnali di allarme precoce.
\item
  \textbf{Cap. 11} --- Predizioni rischiose datate: aggregazione delle
  \(4 \times 3 = 12\) predizioni dei Cap. 7-9 in un \emph{bouquet di
  falsificazione pubblica} della DEQ$^2$, con date, soglie numeriche, fonti
  dati.
\end{itemize}


